\section{Introduction}
\label{sec:intro}

Solar hydrogen is attractive because it stores intermittent sunlight in a
dispatchable chemical bond, yet particulate photocatalysis remains limited by
a chain of losses -- photon harvesting, exciton dissociation, charge
extraction, and surface turnover -- each of which can silently consume an
order of magnitude in efficiency~\cite{tang2008mechanism,ma2022charge,
gu2023recent}. Early oxide studies established that carrier lifetime, not
absorption, is often the controlling variable: hole lifetimes in TiO$_2$
correlate directly with water-splitting ability, and oxygen quantum yields
depend strongly on light intensity~\cite{tang2008mechanism}. Any candidate
material class therefore has to be judged not on a single headline rate but
on how well its structure can be engineered along every link of that chain
~\cite{choi2024understanding,ghosh2020identification}.

Covalent organic frameworks (COFs) enter this competition with an unusual
proposition: fully organic, crystalline, porous semiconductors whose band
structure, exciton physics, and active-site placement are all encoded in
molecular building blocks~\cite{mishra2025covalent,qi2023recent,
lopezmagano2022recent,romero2021light}. The reticular toolbox unites the
tunability of molecular photochemistry with the robustness and long-range
order of extended solids~\cite{guntermann2026recent,khalil2024two,
bagheri2020towards}, and two-dimensional (2D) COFs in particular stack
conjugated sheets into columnar $\pi$-arrays that support exciton migration
and out-of-plane charge transport~\cite{zhou2021peg,wan2020a,jin20192d}.
Since the first demonstration that a hydrazone-linked COF evolves hydrogen
under visible light, the field has expanded through azine, imine,
$\beta$-ketoenamine, triazine, and vinylene linkage generations
~\cite{rodriguezcamargo2026a,haase2017structure,xue2023research,
peng2025sp2}, with recent systems reporting apparent quantum efficiencies
(AQE) above 80\% at 450~nm~\cite{li2022covalent} and mass-normalized rates
approaching 400~mmol\,g$^{-1}$\,h$^{-1}$~\cite{zhao2024nanoscale,gao20242d}.
Designability, however, does not by itself establish efficient charge
separation or scalable operation: strongly bound Frenkel excitons, layer
disorder, and cocatalyst dependence remain recurring bottlenecks
~\cite{bao2025excitonic,ran2024progress,zhang2025dual}.

The literature now contains hundreds of COF photocatalysts but comparatively
few controlled comparisons, and reviews that organize the material by
publication date tend to obscure which structural lever actually moved which
observable~\cite{aslam2026a,you2021recent,yang2020covalent,zhou2023recent}.
This Review therefore organizes evidence by the causal design chain
summarized in Figure~\ref{fig:factorchain}: linkage chemistry $\rightarrow$
conjugation and planarity $\rightarrow$ band gap and band-edge positions
$\rightarrow$ exciton dissociation and charge separation $\rightarrow$
cocatalyst interface $\rightarrow$ HER performance. Each design lever --
donor--acceptor (D--A) backbone engineering, linkage protonation, single-atom
cocatalysts, and Z/S-scheme heterojunctions -- is treated as an intervention
on a specific link, and mechanistic spectroscopy is treated as the
measurement that can distinguish competing explanations
~\cite{ghosh2020identification,choi2024understanding,shen2023efficient}.

The remainder of the Review follows this chain. Section~\ref{sec:levers}
maps the linkage families and the design levers acting on band structure and
exciton physics. Section~\ref{sec:mechanism} reviews the mechanistic
characterization toolbox, from femtosecond transient absorption to operando
spectroscopy. Section~\ref{sec:synthesis} examines synthesis methodology and
crystallinity control, the practical variables that gate every other claim.
Section~\ref{sec:performance} compares reported performance under explicit
protocol variables. Section~\ref{sec:idea} converts the accumulated design
rules into a pre-registered, falsifiable improvement route for the workhorse
$\beta$-ketoenamine COF TpPa-1. Section~\ref{sec:challenges} confronts the
stability, measurement, and reactor questions that stand between
laboratory AQE and solar fuel production, and Section~\ref{sec:conclusion}
concludes.
